\section{周期缺陷约束}
\label{sec:defect-obstruction}

得到有限 survivor 表所用的矩恒等式并不依赖周期很小。下面记录其一般后果，从而把任意
低谱 word 必须满足的局部几何条件，与周期八特有的精确分类区分开来。

\subsection{任意周期 word}

设 $\tau$ 有任意展示周期 $p$，令 $Q_i=\tau_i\tau_{i+1}$，并按
\eqref{eq:defect-statistics} 定义 $d,a,b$。

\begin{theorem}[矩约束]
\label{thm:moment-obstruction}
若 $R_p(\tau)\leq8$，则
\begin{equation}
 d\leq\frac{3p}{4},\qquad 40d+96a+48b\leq42p.
 \label{eq:general-defect-bounds}
\end{equation}
\end{theorem}

\begin{proof}
相位平均恰好筛选净胞元位移为零的闭游走，因此局部展开给出
\[
M_1=4p,\qquad M_2=20p+16d,\qquad
M_3=118p+168d+96a+48b.
\]
谱边假设给出 $M_2\leq8M_1$ 与 $M_3\leq8M_2$，代入化简即得结论。
\end{proof}

\subsection{密度与聚集}

第一个不等式限制缺陷密度；第二个还分别以 $a,b$ 惩罚相邻及距离二的正元对。在固定
$d$ 时，把缺陷聚集只会使必要条件更难满足。这正是平方算子中消去机制在任意周期上的
表现：负 $Q$-位置抑制奇位移耦合，正位置重新激活耦合，而下一阶矩探测短程聚集。

这些不等式仅为必要条件，既不描述所有谱边不超过 $8$ 的长周期 word，也不求出
\eqref{eq:fixed-minimum} 的全局最小值。它们的准确作用是解释短周期候选为何缩减为稀疏
缺陷，以及对径八周期几何为何与较低谱边相容。本文结果不需要更高矩 catalogue。
